% ---------------------------------------------------------------------------
% Theory section. Requires amsmath, amssymb, amsthm and the environments
% theorem, lemma, proposition, corollary, remark (declared in the preamble).
% Proofs are in theory_appendix.tex.
% ---------------------------------------------------------------------------
\section{When can the anomalous direction be identified?}\label{sec:theory}

This section explains the empirical pattern of Sections~\ref{sec:rq3}--\ref{sec:rq5}: label-free rules fail, labelled normal nodes barely help, and a few labelled anomalies are enough. The results are stated for a single relational or attribute statistic on a fixed graph, so they hold without any distributional model of the graph. Proofs are given in Appendix~\ref{app:proofs}.

\subsection{Setting}

Let $G$ be a graph with node set $V$, $|V| = n$, and let $\mathcal{A} \subset V$ be the anomalous nodes, $|\mathcal{A}| = m$, with contamination $\pi = m/n$; $\mathcal{N} = V \setminus \mathcal{A}$ are the normal nodes. A statistic assigns a value $t_i$ to every node. Its \emph{graph-wide rank position} is
\begin{equation}\label{eq:rankpos}
U_i \;=\; \frac{1}{n}\sum_{l \in V}\Big(\mathbb{1}[t_l < t_i] + \tfrac12\,\mathbb{1}[t_l = t_i]\Big) \in (0,1),
\end{equation}
which equals $(r_i + 1)/2$ for the graph-wide signed midrank $r_i$ used by the detector. Write
\begin{equation}
\mu = \frac{1}{m}\sum_{i \in \mathcal{A}} U_i, \qquad
\mathrm{AUC} = \frac{1}{m\,(n-m)}\sum_{i \in \mathcal{A}}\sum_{l \in \mathcal{N}}\Big(\mathbb{1}[t_l < t_i] + \tfrac12\,\mathbb{1}[t_l = t_i]\Big),
\end{equation}
so $\mathrm{AUC}$ is the AUROC obtained by scoring nodes with the upper tail of the statistic. The upper tail is the anomalous direction if $\mathrm{AUC} > \tfrac12$ and the lower tail if $\mathrm{AUC} < \tfrac12$.

\begin{lemma}[Direction is the sign of the mean rank position]\label{lem:sign}
$\mu - \tfrac12 = (1-\pi)\,(\mathrm{AUC} - \tfrac12)$. Hence the upper tail is anomalous if and only if $\mu > \tfrac12$, and the \emph{direction margin}
\[ \Delta = |\mu - \tfrac12| = (1-\pi)\,|\mathrm{AUC} - \tfrac12| \]
measures how strongly the statistic separates anomalies in its anomalous direction.
\end{lemma}

The lemma matters because $\mu$ can be estimated from labelled anomalies alone, using the whole unlabelled graph as reference: no labelled normal node is needed.

\subsection{Unlabelled data do not identify the direction}

A \emph{label-free rule} is any procedure, possibly randomised, whose output depends on the graph, its features and the statistics, but not on any label. All 16 rules of Section~\ref{sec:rq4} are of this kind.

\begin{theorem}[Label-free non-identifiability]\label{thm:nonid}
Let the statistic take distinct values and let $1 \le m < n/2$. Let $\mathcal{A}_{\uparrow}$ be the $m$ nodes with the largest values and $\mathcal{A}_{\downarrow}$ the $m$ nodes with the smallest. Labelling $\mathcal{A}_{\uparrow}$ as anomalous gives $\mathrm{AUC} = 1$; labelling $\mathcal{A}_{\downarrow}$ gives $\mathrm{AUC} = 0$. Both labellings have the same contamination and leave the input of every label-free rule unchanged, so every label-free rule chooses the wrong signed direction (upper or lower) with probability at least $\tfrac12$ under one of them.
\end{theorem}

Theorem~\ref{thm:nonid} is elementary, but it fixes what a label-free rule can and cannot do. Skewness, consensus and similar rules succeed only when the unknown labelling mechanism happens to place anomalies in the tail the rule favours. That is an assumption about how anomalies arise, not something the data can confirm. The benchmarks contain both mechanisms: injected clique anomalies sit in the upper tail of triangle statistics, while organic camouflaged anomalies sit in the lower tail of neighbour disagreement (Section~\ref{sec:rq3}). This is why no fixed label-free rule transferred.

\begin{remark}[Scope of Theorem~\ref{thm:nonid}]\label{rem:scope}
The theorem concerns the \emph{signed} direction. It does not say that unlabelled data are useless for detection: two-sided evidence $|r|$ needs no labels and scores both tails. Its cost is that half of its evidence lies in the tail that contains no anomalies, so it ranks the anomalies below every extreme normal node in the other tail. In the earlier study the all-two-sided rule was a reasonable label-free default (mean AUROC 0.66--0.68 across the three splits) but was beaten by correctly signed statistics wherever the direction was stable. FDC therefore keeps the two-sided option and lets the labelled anomalies choose among the three.
\end{remark}

\subsection{Labelled normal nodes are weak evidence}

Labelled normal nodes add information, but very little when anomalies are rare. Suppose $m'$ normal nodes are drawn uniformly without replacement from $\mathcal{N}$ and revealed.

\begin{theorem}[Normal shots]\label{thm:normal}
In the setting of Theorem~\ref{thm:nonid}, assume $n \ge 2m + m'$. Every rule that uses the graph and $m'$ labelled normal nodes chooses the wrong direction, under one of the two labellings, with probability at least
\[
\tfrac12\,q, \qquad
q = \prod_{i=0}^{m'-1}\frac{n-2m-i}{\,n-m-i\,} \;\ge\; \Big(1 - \frac{m}{\,n-m-m'+1\,}\Big)^{m'} \;\approx\; \exp\!\Big(-\frac{m'\pi}{1-\pi}\Big).
\]
\end{theorem}

To push this worst-case error below $\varepsilon$, roughly $m' \ge \frac{1-\pi}{\pi}\ln\frac{1}{2\varepsilon}$ normal nodes are required. With $\pi = 0.05$ and $\varepsilon = 0.1$ this is about 31 labelled normal nodes for a single statistic. With the $m' = 10$ normal shots used by ARC and AnomalyGFM, the bound is between $0.15$ and $0.39$ on the 14 of the 16 development graphs with $\pi < 0.11$. This agrees with the near-chance direction recovery from normal shots reported in Section~\ref{sec:rq5}. The reason is an asymmetry: when anomalies are rare, normal nodes look like the graph as a whole, so they say little about which tail the few anomalies occupy.

\subsection{A few labelled anomalies identify the direction}

Now suppose $k$ anomalies are drawn uniformly without replacement from $\mathcal{A}$ and revealed. The calibration rule estimates $\mu$ by the mean rank position of the shots,
\[
\hat\mu = \frac1k \sum_{s=1}^{k} U_{a_s},
\]
and chooses the upper tail if $\hat\mu > \tfrac12$ and the lower tail otherwise. The mean of $\hat\mu$ is $\mu$, and each $U_{a_s}$ lies in $(0,1)$.

\begin{theorem}[Anomaly shots]\label{thm:anomaly}
Let $\Delta > 0$ be the direction margin of Lemma~\ref{lem:sign} and $\sigma^2 = \frac1m \sum_{i \in \mathcal{A}} (U_i - \mu)^2$ the spread of the anomalies' rank positions. The probability that the calibration rule chooses the wrong direction satisfies
\[
\Pr[\text{wrong direction}] \;\le\; \min\left\{ \exp\!\big(-2k\Delta^2\big),\;
\exp\!\left(-\frac{k\Delta^2}{2\sigma^2 + 2\Delta/3}\right)\right\}.
\]
\end{theorem}

Theorem~\ref{thm:anomaly} covers the choice between the two signed directions. FDC chooses among three candidates, the two signed tails and the two-sided evidence, by the largest shot mean. The same argument applies to this three-way choice.

\begin{corollary}[Three-way selection]\label{cor:three}
Let $\mu^{\uparrow}, \mu^{\downarrow}, \mu^{\leftrightarrow}$ be the mean upper, lower and two-sided evidence of the anomalies, let $t^{*}$ be the candidate with the largest mean, and let $g_t = \mu^{t^*} - \mu^{t} > 0$ be its gap to each other candidate $t$. The probability that the largest shot mean belongs to a candidate other than $t^*$ is at most $\sum_{t \ne t^*} \exp(-k g_t^2/2)$. For the two signed tails $g = 2\Delta$, which recovers the Hoeffding term of Theorem~\ref{thm:anomaly}.
\end{corollary}

\begin{corollary}[How many anomalies to label]\label{cor:k}
If $J$ statistics all have margin at least $\Delta_{\min}$, all their directions are correct with probability at least $1-\varepsilon$ once
\[
k \;\ge\; \frac{\ln(J/\varepsilon)}{2\,\Delta_{\min}^{2}}.
\]
\end{corollary}

The contrast with Theorem~\ref{thm:normal} is sharp. In the two extreme labellings of Theorem~\ref{thm:nonid}, a single labelled anomaly reveals the direction exactly, while $m'$ normal nodes leave it undetermined with probability close to $\exp(-m'\pi/(1-\pi))$. In general the error with anomaly shots decays exponentially in $k$ at a rate set by $\Delta$, which depends on the contamination only through the factor $1-\pi$. With labelled normal nodes, the worst-case rate is throttled by $\pi$ itself.

\subsection{Choosing which statistics to trust}

With few shots, some statistics look separating by chance. The method therefore weights statistic $j$ by $w_j = \max(0,\ \hat s_j - \tau_k)$. Here $\hat s_j$ is the estimated separation of its chosen tail and $\tau_k = 1/\sqrt{12k}$.

\begin{proposition}[Null behaviour of the evidence]\label{prop:tau}
Suppose the $k$ shots are a uniform random sample of nodes with respect to statistic $j$ (no anomaly signal).
\begin{enumerate}
\item[(a)] For the upper and lower evidence, with or without ties, the shot mean has expectation exactly $\tfrac12$ and standard deviation at most $1/\sqrt{12k} = \tau_k$.
\item[(b)] For the two-sided evidence, the shot mean has expectation $m_j = \frac1n\sum_i |\bar r_{ij}| \le \tfrac12$, with equality up to $O(1/n^2)$ when the values are distinct, and variance $\sigma_{j}^2/k$, where $\sigma_j^2 = \frac1n\sum_i(|\bar r_{ij}| - m_j)^2 \le \tfrac14$. With ties, $m_j$ can be well below $\tfrac12$ and $\sigma_j^2$ can exceed $\tfrac1{12}$.
\item[(c)] If instead statistic $j$ has margin $\Delta_j > \tau_k$, it keeps a positive weight with probability at least $1 - \exp\!\big(-2k(\Delta_j - \tau_k)^2\big)$.
\end{enumerate}
\end{proposition}

For the signed evidence, $\tau_k$ is therefore an upper bound on one null standard deviation of the shot mean, and statistics whose apparent separation lies within it are switched off. For the two-sided evidence the method centres at $\tfrac12 \ge m_j$, which makes retention harder when ties lower $m_j$, but the null spread can exceed $\tau_k$. On our graphs this matters little: across the 143 non-empty statistic--graph pairs of the development set, $\sigma_j^2 \le 0.0843$, at most $1.2\%$ above $\tfrac1{12}$, while heavily tied statistics have $m_j$ as low as $0.28$ (and $0$ for the few statistics that are constant on a graph). The threshold is a per-statistic noise floor, not a guarantee on the number of statistics retained by chance across all $9+2p$ statistics and three candidate directions; Section~\ref{sec:null} measures that number directly with random shots. Finally, the proposition explains why few shots call for caution with many statistics (for $k=1$, a signed statistic with distinct values and no signal passes the threshold with probability $0.42$), but it does not single out $k = 3$ as the point at which to add attribute statistics. That gate was chosen in the last development round (Section~\ref{sec:setup}); Section~\ref{sec:sens} reports the alternatives.

\subsection{Empirical check}

Table~\ref{tab:bound} compares the observed direction error of the calibration rule with Theorem~\ref{thm:anomaly}, using the exact $\Delta$ and $\sigma^2$ of each statistic computed from all labels. It pools 9 relational statistics, 16 development graphs and 5 seeds. The observed error never exceeds either bound. Both bounds are conservative: for margins above $0.2$, no error was observed at $k = 10$, while Hoeffding's bound allows $35\%$ and Bernstein's $16\%$. The variance-adaptive Bernstein bound is the tighter of the two throughout. Statistics with margin below $0.05$ are close to a coin flip at every $k$, as the theory predicts; the weight threshold of Proposition~\ref{prop:tau} is designed to discard exactly these.

\begin{table}[t]
\caption{Observed direction error of the anomaly-shot calibration rule against the bounds of Theorem~\ref{thm:anomaly}. Rows pool statistics by their margin $\Delta$ into five intervals, all of which are reported for every $k$; $n$ is the number of (statistic, graph, seed) cases. Development graphs only.}\label{tab:bound}
\centering\small
\begin{tabular}{@{}cl rrrr@{}}
\toprule
$k$ & $\Delta$ & $n$ & Observed & Hoeffding & Bernstein \\
\midrule
1 & $(0.00, 0.05]$ & 172 & 0.494 & 0.999 & 0.997 \\
 & $(0.05, 0.10]$ & 123 & 0.374 & 0.988 & 0.973 \\
 & $(0.10, 0.20]$ & 155 & 0.361 & 0.956 & 0.915 \\
 & $(0.20, 0.30]$ & 113 & 0.177 & 0.898 & 0.828 \\
 & $(0.30, 0.50]$ & 99 & 0.030 & 0.760 & 0.630 \\
\midrule
3 & $(0.00, 0.05]$ & 182 & 0.467 & 0.996 & 0.989 \\
 & $(0.05, 0.10]$ & 131 & 0.229 & 0.964 & 0.920 \\
 & $(0.10, 0.20]$ & 158 & 0.152 & 0.876 & 0.770 \\
 & $(0.20, 0.30]$ & 120 & 0.092 & 0.725 & 0.571 \\
 & $(0.30, 0.50]$ & 104 & 0.019 & 0.444 & 0.258 \\
\midrule
5 & $(0.00, 0.05]$ & 178 & 0.421 & 0.994 & 0.982 \\
 & $(0.05, 0.10]$ & 116 & 0.353 & 0.939 & 0.865 \\
 & $(0.10, 0.20]$ & 138 & 0.188 & 0.795 & 0.647 \\
 & $(0.20, 0.30]$ & 120 & 0.050 & 0.587 & 0.396 \\
 & $(0.30, 0.50]$ & 100 & 0.000 & 0.259 & 0.108 \\
\midrule
10 & $(0.00, 0.05]$ & 178 & 0.478 & 0.987 & 0.965 \\
 & $(0.05, 0.10]$ & 117 & 0.205 & 0.882 & 0.751 \\
 & $(0.10, 0.20]$ & 139 & 0.072 & 0.639 & 0.429 \\
 & $(0.20, 0.30]$ & 121 & 0.000 & 0.351 & 0.163 \\
 & $(0.30, 0.50]$ & 100 & 0.000 & 0.073 & 0.015 \\
\bottomrule
\end{tabular}
\end{table}
